\section{Incident Summary}

\subsection{Observed Event}
This report is intended for structured failure investigation, field return
analysis, unexpected lab events, or systematic fault diagnosis. It should make
the symptom, context, and containment state clear before hypotheses start to
compete with one another.

\dangerbox{
Do not open the root-cause section with a guess. Start with the symptom,
environment, and impact. Premature certainty is one of the fastest ways to
damage a good investigation.
}

\subsection{Event Timeline}
\begin{table}[H]
    \centering
    \small
    \begin{tabular}{p{2.4cm} p{3.0cm} p{5.5cm}}
        \toprule
        \textbf{Timestamp} & \textbf{Observed By} & \textbf{Event} \\
        \midrule
        09:15 & Test Engineer & Unit powered successfully and reached nominal operation. \\
        09:23 & Test Engineer & Output collapsed during dynamic load transition. \\
        09:25 & Lab Lead & Containment applied and additional stress testing stopped. \\
        10:10 & Analysis Team & Preliminary evidence collection started. \\
        \bottomrule
    \end{tabular}
    \caption{Example incident timeline}
    \label{tab:incident-timeline}
\end{table}

\subsection{Impact Statement}
\begin{itemize}
    \item describe the operational impact on the system or customer-visible behavior;
    \item record whether hardware damage, safety risk, or data corruption was observed;
    \item state whether the issue is isolated, intermittent, or reproducible.
\end{itemize}
